% ============================================================
%  Cognitive Canvas: An AI-Powered Local-First Spatial
%  Journaling Framework
%  Single-Column Preprint Format (arXiv / ACM-style)
% ============================================================
\documentclass[12pt, letterpaper]{article}

% ---- Core layout -------------------------------------------
\usepackage[
  top=1.25in,
  bottom=1.25in,
  left=1.25in,
  right=1.25in
]{geometry}

% ---- Typography --------------------------------------------
\usepackage[T1]{fontenc}
\usepackage[utf8]{inputenc}

\usepackage[expansion=false]{microtype}
\usepackage{setspace}
\onehalfspacing

% ---- Math --------------------------------------------------
\usepackage{amsmath,amssymb,amsfonts}

% ---- Graphics & color --------------------------------------
\usepackage{graphicx}
\usepackage[dvipsnames,svgnames,table]{xcolor}

% ---- Tables ------------------------------------------------
\usepackage{booktabs}
\usepackage{array}
\usepackage{tabularx}
\usepackage{multirow}

% ---- Lists -------------------------------------------------
\usepackage{enumitem}
\setlist[enumerate]{leftmargin=1.8em, itemsep=3pt, parsep=0pt}
\setlist[itemize]{leftmargin=1.8em,  itemsep=3pt, parsep=0pt}

% ---- Section styling ---------------------------------------
\usepackage{titlesec}
\titleformat{\section}
  {\large\bfseries\scshape}{\Roman{section}.}{0.6em}{}
  [\vspace{-2pt}\rule{\linewidth}{0.4pt}\vspace{4pt}]
\titleformat{\subsection}
  {\normalsize\bfseries\itshape}{\Alph{subsection}.}{0.5em}{}
\titleformat{\subsubsection}[runin]
  {\normalsize\bfseries}{}{0em}{}[.\enspace]
\titlespacing{\section}{0pt}{18pt}{6pt}
\titlespacing{\subsection}{0pt}{12pt}{4pt}
\titlespacing{\subsubsection}{0pt}{8pt}{4pt}

% ---- Header / footer ---------------------------------------
\usepackage{fancyhdr}
\pagestyle{fancy}
\fancyhf{}
\fancyhead[L]{\small\textit{Cognitive Canvas — Preprint}}
\fancyhead[R]{\small\textit{Siva Prakash R K et al.}}
\fancyfoot[C]{\thepage}
\renewcommand{\headrulewidth}{0.4pt}

% ---- Captions ----------------------------------------------
\usepackage[font=small, labelfont=bf, labelsep=period,
            justification=justified]{caption}

% ---- Hyperlinks (colored, website-style) -------------------
\usepackage{hyperref}
\hypersetup{
  colorlinks  = true,
  linkcolor   = MidnightBlue,
  citecolor   = ForestGreen,
  urlcolor    = RoyalBlue,
  pdftitle    = {Cognitive Canvas: An AI-Powered Local-First Spatial Journaling Framework},
  pdfauthor   = {Siva Prakash R K et al.},
  pdfsubject  = {Spatial Journaling, Privacy-Preserving AI, NLP},
  pdfkeywords = {Spatial Journaling, Privacy AI, RAG, NLP, Knowledge Graph},
  bookmarks   = true,
  pdfpagemode = UseOutlines
}

% ---- Abstract box ------------------------------------------
\usepackage{mdframed}
\mdfdefinestyle{abstractstyle}{
  linecolor    = gray!40,
  linewidth    = 0.8pt,
  leftmargin   = 0pt,
  rightmargin  = 0pt,
  innerleftmargin  = 14pt,
  innerrightmargin = 14pt,
  innertopmargin   = 12pt,
  innerbottommargin= 12pt,
  backgroundcolor  = gray!5,
  skipabove    = 12pt,
  skipbelow    = 12pt
}

% ---- Cite --------------------------------------------------
\usepackage{cite}

% ---- URL ---------------------------------------------------
\usepackage{url}
\urlstyle{same}

% ---- Macros ------------------------------------------------
\newcommand{\CC}{\textit{Cognitive Canvas}}
\newcommand{\sys}{\textit{CC}}
\newcommand{\eg}{e.g.,\xspace}
\newcommand{\ie}{i.e.,\xspace}

% ---- Colored rule for title block --------------------------
\newcommand{\mytitlerule}{%
  {\color{MidnightBlue!60}\rule{\linewidth}{1.5pt}}%
}

% ============================================================
\begin{document}

% ---- Remove page number on title page ----------------------
\thispagestyle{empty}

% ============================================================
% TITLE BLOCK
% ============================================================

\vspace*{0.3in}
\mytitlerule
\vspace{10pt}

{\LARGE\bfseries Cognitive Canvas: An AI-Powered Local-First\\[4pt]
Spatial Journaling Framework}

\vspace{14pt}

{\large
Siva Prakash R K\textsuperscript{*},\;
Kiruba Sri G\textsuperscript{*},\;
Priyangga T\textsuperscript{*},\;
Santhosh S\textsuperscript{*},\;
Rishi Shankar M\textsuperscript{*}
}

\vspace{6pt}

{\normalsize\textit{
\textsuperscript{*}Department of Computer Science and Engineering,
[Insert College Name], Salem, Tamil Nadu, India
}}

\vspace{4pt}

{\small
\texttt{\{[insert.email1],\,[insert.email2],\,[insert.email3]\}@institution.edu}
}

\vspace{10pt}
\mytitlerule
\vspace{18pt}

% ============================================================
% ABSTRACT
% ============================================================

\begin{mdframed}[style=abstractstyle]
\noindent\textbf{\textsc{Abstract} ---}
Contemporary digital journaling applications fail to surface
longitudinal psychological patterns, and cloud-dependent AI
augmentations introduce acute privacy risks incompatible with
the intimate nature of personal reflection data.
This paper presents \CC{}, a local-first, privacy-preserving
spatial journaling framework that transforms free-form textual
entries into an interactive, node-based knowledge graph
entirely on consumer-grade hardware.
The system encodes each journal entry into dense semantic
vectors using a locally quantised variant of RoBERTa, stores
and retrieves those vectors at sub-second latency via the
Facebook AI Similarity Search (FAISS) library, and interprets
user queries through an on-device Rasa Natural Language
Understanding (NLU) pipeline---ensuring that no user data
traverses a network boundary at any stage.
Retrieval-Augmented Generation (RAG) over the personal vector
store enables contextually grounded, temporally aware responses
without cloud inference.
A pilot evaluation on consumer-grade laptop hardware (12th Gen
Intel Core i5-12450HX, 16\,GB RAM, no discrete GPU) demonstrates
end-to-end query latency below 1.2\,s and vector-insert
throughput exceeding 800 entries per minute.
Semantic link precision, assessed across 35 manually annotated
entry pairs, reached 80\%, compared with 60\% under legacy
static-threshold approaches.
\CC{} establishes a reproducible architectural baseline for
privacy-first, spatially rendered cognitive self-reflection
tools.

\vspace{8pt}
\noindent\textbf{Keywords:}
Spatial Journaling,\;
Privacy-Preserving AI,\;
Retrieval-Augmented Generation,\;
Natural Language Processing,\;
Local-First Architecture,\;
Knowledge Graph,\;
Cognitive Drift Detection
\end{mdframed}

\newpage

% ============================================================
\section{Introduction}
% ============================================================

Personal journaling is among the most empirically validated
interventions for emotional regulation, stress reduction, and
metacognitive development~\cite{pennebaker1997,smyth1998}.
Despite this, the transition of journaling into the digital
domain has largely preserved its most limiting structural
property: strict chronological linearity.
Entries accumulate as flat text sequences, burying temporally
distant but thematically proximate reflections beneath layers
of irrelevant content and rendering long-range psychological
pattern detection practically infeasible for both user and
system.

The advent of large language models (LLMs) and semantic search
has prompted a wave of ``AI-assisted journaling'' products that
promise to surface such patterns
automatically~\cite{xu2023mindfuldiary}.
However, practically all production deployments in this category
transmit entry content to remote inference endpoints, exposing
intimate self-reflection data to third-party servers,
algorithmic surveillance, and the ever-present risk of data
breach~\cite{feretzakis2024}.
For a category of data whose therapeutic value depends
critically on uninhibited, private disclosure, this
architectural decision constitutes a fundamental contradiction.

\subsection*{Gaps in Current Applications}

A systematic examination of existing AI journaling tools
reveals four persistent deficiencies:

\begin{enumerate}

  \item \textbf{Cloud-mandatory inference.}
    Tools such as MindfulDiary~\cite{xu2023mindfuldiary}
    transmit raw entry text to remote LLM APIs.
    Users possess no verifiable guarantee of data isolation,
    retention policy compliance, or inference confidentiality.

  \item \textbf{Static semantic thresholding.}
    Systems that do perform local similarity computations
    typically apply fixed cosine-distance cutoffs (commonly
    $\tau = 0.78$~\cite{tauscher2023}), which fail to adapt
    to the idiosyncratic vocabulary distributions of individual
    users, producing high false-positive connection rates in
    verbose writers and sparse graphs in terse ones.

  \item \textbf{Absence of temporal topology.}
    Existing knowledge-graph journaling prototypes represent
    entries as nodes in undirected, force-directed graphs that
    encode thematic similarity but discard temporal
    sequence---erasing the directional arc of cognitive
    evolution that is arguably the most diagnostically relevant
    signal.

  \item \textbf{Clinical framing without clinical validation.}
    Several systems label inferred emotional states with
    clinical terminology (e.g., ``anxiety'', ``depression
    risk'') derived from classifiers trained on datasets
    exhibiting severe population and annotation
    biases~\cite{sage2025}, creating potential for harm without
    providing actionable clinical value.

\end{enumerate}

\CC{} is designed to close all four gaps simultaneously.
The framework runs its entire inference pipeline on local
consumer hardware, employs a user-adaptive semantic
thresholding protocol (Adaptive Min-P), renders entries on a
deterministic Rigid Narrative Timeline that preserves
chronological order while overlaying AI-inferred thematic arcs,
and deliberately maps a non-pathological \emph{Cognitive Drift}
metric in place of clinical labels.

The remainder of this paper is structured as follows.
Section~II surveys related work in AI-assisted journaling and
local RAG pipelines.
Section~III formalises the problem and motivates the
architectural choices.
Section~IV presents the staged system architecture and
enumerates the framework's innovative contributions.
Section~V reports pilot performance results and discusses
expected system behaviour.
Section~VI concludes and outlines directions for future work.

% ============================================================
\section{Related Work}
% ============================================================

\subsection{AI-Augmented Journaling Systems}

Early computational approaches to journaling analysis focused
on lexicon-based affect scoring
(LIWC~\cite{pennebaker2003liwc}, VADER), producing aggregate
sentiment statistics rather than relational structure.
More recent systems leverage transformer-based language models
to extract deeper semantic representations.
MindfulDiary~\cite{xu2023mindfuldiary} employs GPT-4 as a
conversational journaling partner, demonstrating measurable
improvements in user reflection depth; however, the system
requires a persistent internet connection and routes all data
through OpenAI's cloud infrastructure.
MindScape~\cite{mindscape} extends journaling with passive
physiological sensing, but its reliance on invasive sensor
data introduces a distinct privacy surface.
C-Journal~\cite{cjournal2024} adopts a graph-based
representation of journal entries, yet applies a static
similarity threshold ($\tau = 0.78$) and does not support
local execution.

\subsection{Retrieval-Augmented Generation}

The RAG paradigm~\cite{lewis2020rag}, which conditions
generative model outputs on documents retrieved from a
personal or domain-specific corpus, has emerged as the
dominant architecture for grounding LLM responses in
verifiable, user-specific context.
Canonical RAG pipelines encode a document corpus into dense
vector embeddings, index those vectors in an approximate
nearest-neighbour (ANN) store, and retrieve the top-$k$
most semantically similar passages at inference
time~\cite{karpukhin2020dpr}.
FAISS~\cite{johnson2021faiss} provides a highly optimised,
CPU-amenable ANN index that has been successfully deployed in
resource-constrained edge settings.
Recent work on Small Language Models (SLMs) and quantised
transformer inference demonstrates that sub-billion parameter
models can achieve competitive semantic encoding quality on
consumer CPUs at acceptable
latency~\cite{dettmers2022llmint8}.

\subsection{Critical Gap}

Despite the theoretical compatibility of RAG with local
execution, no published system applies the full RAG
pipeline---embedding, retrieval, and generation---to personal
journaling in a fully local, zero-network configuration, with
user-adaptive retrieval thresholds and temporally structured
visualisation.
The present work addresses this gap directly.

% ============================================================
\section{Motivation, Problem Statement, and Proposed Methodology}
% ============================================================

\subsection{Motivation}

The motivating observation is straightforward: the therapeutic
efficacy of journaling depends on candid, uninhibited
disclosure, yet every major AI journaling product requires the
user to transmit that candid disclosure to a remote server.
This contradiction is not merely a privacy preference---it is a
structural barrier to adoption by any user population with
legitimate confidentiality concerns (mental health patients,
domestic abuse survivors, corporate employees subject to
non-disclosure obligations, minors).
A privacy-preserving alternative that does not sacrifice
semantic intelligence is therefore not a niche optimisation
but a prerequisite for equitable access.

\subsection{Problem Formulation}

Let $\mathcal{E} = \{e_1, e_2, \ldots, e_n\}$ denote a
chronologically ordered corpus of journal entries authored by
a single user, where each entry $e_i$ is a natural language
text document associated with a timestamp $t_i$.
The system must:

\begin{enumerate}
  \item Compute a dense vector representation
    $\mathbf{v}_i \in \mathbb{R}^d$ for each $e_i$
    using only local computation.

  \item Maintain a persistent, queryable index
    $\mathcal{I}$ over $\{\mathbf{v}_i\}$ supporting
    sub-second approximate nearest-neighbour retrieval.

  \item Given a natural language query $q$ from the user,
    classify the intent $\hat{\iota} = f(q)$, retrieve the
    top-$k$ most relevant entries
    $R_q \subset \mathcal{E}$, and generate a contextually
    grounded response $r$ conditioned on $R_q$---all without
    network egress.

  \item Construct and maintain a directed graph
    $G = (V, E)$ where $V = \mathcal{E}$ and edges $E$
    encode semantic, contextual, and compensatory
    relationships according to user-adaptive thresholds.
\end{enumerate}

Formally, the privacy constraint is:
$\forall\, e_i \in \mathcal{E},\ \forall\, q:\
\nexists$ network transmission of $e_i$, $q$,
$\mathbf{v}_i$, or any derivative thereof.

\subsection{Proposed Methodology}

\CC{} satisfies the above formulation through three tightly
integrated local components.

\subsubsection{RoBERTa Embedding Engine}
Each entry is encoded by a locally deployed, INT8-quantised
\texttt{roberta-base} variant fine-tuned for semantic textual
similarity~\cite{liu2019roberta}.
Quantisation reduces peak RAM consumption from approximately
500\,MB (FP32) to 130\,MB (INT8), enabling parallel execution
alongside the inference runtime on a 16\,GB system.
The resulting 768-dimensional embedding $\mathbf{v}_i$ is
L2-normalised prior to indexing, converting cosine similarity
queries into equivalent inner-product searches supported
natively by FAISS.

\subsubsection{FAISS Vector Store}
Normalised embeddings are inserted into a FAISS
\texttt{IndexFlatIP} (exact inner-product) index persisted to
a local \texttt{sqlite-vec} database.
For corpora exceeding 10,000 entries, the index transitions
automatically to an \texttt{IndexIVFFlat} configuration with
$n_{\mathrm{list}} = \lfloor\sqrt{N}\rfloor$ Voronoi cells,
maintaining sub-100\,ms retrieval latency at
scale~\cite{johnson2021faiss}.
All index files reside in the user's local application
directory; no synchronisation service is invoked.

\subsubsection{Rasa NLU Intent Pipeline}
User queries are processed by an on-device Rasa NLU pipeline
comprising a SpaCy tokeniser, a CountVectors featuriser, a
DIETClassifier~\cite{bunk2020diet}, and an EntityRecogniser.
The pipeline is trained on a domain-specific intent taxonomy
(Table~\ref{tab:intents}) and executes entirely in-process,
adding fewer than 80\,ms of overhead per query on the
evaluation hardware.

\begin{table}[htbp]
\centering
\caption{Rasa NLU Intent Taxonomy}
\label{tab:intents}
\begin{tabular}{ll}
\toprule
\textbf{Intent Class} & \textbf{Example Utterance} \\
\midrule
\texttt{retrieve\_theme}   & ``Show me entries about anxiety'' \\
\texttt{temporal\_query}   & ``What did I write last week?'' \\
\texttt{emotion\_summary}  & ``How have I been feeling lately?'' \\
\texttt{drift\_report}     & ``Am I thinking about the same things?'' \\
\texttt{add\_entry}        & ``New entry: Today was difficult.'' \\
\texttt{visualise\_graph}  & ``Show me my thought map'' \\
\texttt{rumination\_check} & ``Have I been repeating myself?'' \\
\bottomrule
\end{tabular}
\end{table}

% ============================================================
\section{System Architecture and Innovative Contributions}
% ============================================================

\subsection{Architectural Overview}

\CC{} adopts a layered client-server architecture confined
entirely to the local device.
A FastAPI backend exposes a REST interface consumed by an
HTML/D3.js frontend; both processes bind exclusively to
\texttt{localhost} and are inaccessible from external networks.
The pipeline executes across six discrete stages described
in the subsections below.

\subsection{Stage 1 --- Input Processing and Intent Classification}

Raw text submitted through the journaling interface is first
passed to the Rasa NLU pipeline, which returns a classified
intent $\hat{\iota}$ and a set of extracted entities (date
references, emotion keywords, topic labels).
For \texttt{add\_entry} intents, the preprocessor performs
Unicode normalisation (NFC), punctuation standardisation, and
length-adaptive sentence segmentation before forwarding the
cleaned text to Stage~2.
For retrieval intents, the raw query is forwarded directly to
Stage~3 following entity extraction.

\subsection{Stage 2 --- Vector Encoding}

The sanitised entry text is encoded by the quantised RoBERTa
model.
The forward pass produces a 768-dimensional hidden-state
tensor; the \texttt{[CLS]} token representation is extracted
as the sentence embedding $\mathbf{v}_i$, then L2-normalised:
\begin{equation}
\hat{\mathbf{v}}_i = \frac{\mathbf{v}_i}{\|\mathbf{v}_i\|_2}.
\label{eq:normalise}
\end{equation}
Embedding generation requires approximately 180--220\,ms on
the evaluation hardware (Intel Core i5-12450HX, 16\,GB RAM,
no GPU acceleration).

\subsection{Stage 3 --- Adaptive Thresholding and Graph Linking}

Rather than applying a fixed cosine threshold $\tau$, \CC{}
computes a user-adaptive threshold from the statistical
distribution of the personal embedding corpus:
\begin{equation}
T_{\mathrm{adaptive}} = \mu_{\mathrm{user}} + \sigma_{\mathrm{user}},
\label{eq:adaptive}
\end{equation}
where $\mu_{\mathrm{user}}$ and $\sigma_{\mathrm{user}}$
denote the mean and standard deviation of pairwise cosine
similarities across all entries authored by the current user.
$T_{\mathrm{adaptive}}$ is recomputed incrementally after each
new insertion using Welford's online
algorithm~\cite{welford1962}, adding $O(1)$ overhead per
insert.

The Tripartite Linking Logic evaluates three edge categories:

\begin{itemize}
  \item \textbf{Semantic Chain Links} (Amber):
    For incoming entry $e_n$, the system walks backwards
    through the entry history in reverse chronological order,
    computing cosine similarity against each predecessor until
    the first predecessor whose similarity exceeds
    $T_{\mathrm{adaptive}}$ is identified.
    A single directed semantic edge is inserted; search
    terminates immediately, ensuring at most one semantic
    predecessor per node.

  \item \textbf{Contextual Chain Links} (Blue):
    A directed edge is created from $e_n$ to the most recent
    prior entry sharing at least one explicit user-assigned
    tag, producing tag-based narrative threads rather than
    tag webs.

  \item \textbf{Compensatory Links} (Red):
    Entries exhibiting strongly opposing emotional
    valence---classified by the \texttt{deberta-v3-base}
    zero-shot classifier~\cite{he2021deberta}---with cosine
    similarity $\geq 0.70$ within a 7-day temporal window
    receive a compensatory edge, surfacing psychologically
    significant oscillation patterns.
\end{itemize}

\subsection{Stage 4 --- Cognitive Drift Computation}

The system computes a \emph{Cognitive Drift} score $CD(W)$
for a sliding temporal window $W$:
\begin{equation}
CD(W) = 1 - \cos\!\left(\boldsymbol{\mu}_W,\,
  \boldsymbol{\mu}_{\mathrm{global}}\right),
\label{eq:drift}
\end{equation}
where $\boldsymbol{\mu}_W$ is the centroid of embeddings
within $W$ and $\boldsymbol{\mu}_{\mathrm{global}}$ is the
global embedding centroid over all entries.
$CD(W) \in [0,2]$, with values approaching 0 indicating
topical consistency and values approaching 1 indicating
maximal semantic divergence from the historical mean.
The metric is intentionally non-pathological: it quantifies
directional cognitive change without labelling it as
distorted or disordered.

\subsection{Stage 5 --- Rumination Guard}

The Rumination Guard monitors for sustained semantic loops
co-occurring with negative affect.
Given the current entry $e_n$, the mean embedding of the
$N=5$ most recent entries $\bar{\mathbf{v}}_{n-5:n-1}$, a
similarity threshold $T_{\mathrm{sim}} = 0.88$, and a
negative-valence probability threshold
$P_{\mathrm{neg}} = 0.70$, the guard increments a counter
$k$:
\begin{equation}
k \leftarrow
\begin{cases}
k + 1 & \text{if } \cos(\hat{\mathbf{v}}_n,\, \bar{\mathbf{v}})
  > T_{\mathrm{sim}}\ \text{and}\ P_{\mathrm{neg}} > 0.70 \\
0     & \text{otherwise.}
\end{cases}
\label{eq:rumination}
\end{equation}
When $k \geq K = 3$ consecutive triggers accumulate, the
frontend presents a non-clinical wellness intervention modal
offering structured cognitive reappraisal
prompts~\cite{greenberger1995}, then resets $k$ to 0.
The $K=3$ requirement prevents isolated negative entries from
triggering spurious interventions.

\subsection{Stage 6 --- Rigid Narrative Timeline Visualisation}

Retrieved entries and graph topology are transmitted to the
D3.js frontend via the \texttt{GET /get\_graph\_data}
endpoint.
Unlike conventional force-directed graph layouts, \CC{}
implements a \emph{Rigid Narrative Timeline}: nodes are
positioned deterministically on a horizontal chronological
axis, with newer entries displaced rightward.
AI-inferred arcs are rendered as alternating quadratic
B\'{e}zier curves (alternating above and below the axis per
successive link), preventing visual occlusion between
co-incident edges.
Client-side Transitive Reduction~\cite{aho1972} suppresses
redundant arcs, ensuring the displayed graph encodes only
non-derivable relationships.

\subsection{Human-in-the-Loop Link Verification}

All AI-proposed edges are presented to the user via the
Floating Command Center interface before database persistence.
The two-stage API enforces this protocol structurally:

\begin{itemize}
  \item \texttt{POST /analyze\_thought} --- computes and
    returns candidate links without writing to the database.
  \item \texttt{POST /save\_thought} --- persists the entry
    and only those links in the \texttt{approved\_links}
    array provided by the frontend, populated exclusively by
    explicit user confirmation.
\end{itemize}

Chronological links are the sole exception; they are
auto-verified as factual structural edges.
This design preserves human interpretive authority over the
cognitive map~\cite{amershi2014,horvitz1999}.

\subsection{Innovative Contributions Summary}

Table~\ref{tab:contributions} positions \CC{} relative to
the closest published alternatives across seven architectural
dimensions.

\begin{table}[htbp]
\centering
\caption{Feature Comparison: \CC{} vs.\ Closest Published
Alternatives. \checkmark\,=\,fully supported;
$\times$\,=\,not supported; $\sim$\,=\,partial.}
\label{tab:contributions}
\setlength{\tabcolsep}{7pt}
\begin{tabular}{lccccc}
\toprule
\textbf{Feature} &
  \textbf{\CC{}} &
  \textbf{\cite{xu2023mindfuldiary}} &
  \textbf{\cite{cjournal2024}} &
  \textbf{\cite{mindscape}} &
  \textbf{\cite{johnson2021faiss}} \\
\midrule
Fully local inference         & \checkmark & $\times$   & $\times$  & $\times$  & $\sim$ \\
Adaptive sim.\ threshold      & \checkmark & $\times$   & $\times$  & $\times$  & $\times$ \\
Temporal graph topology       & \checkmark & $\times$   & $\sim$    & $\times$  & $\times$ \\
HITL link verification        & \checkmark & $\times$   & $\times$  & $\times$  & $\times$ \\
Rumination detection          & \checkmark & $\sim$     & $\times$  & $\times$  & $\times$ \\
Non-pathological drift metric & \checkmark & $\times$   & $\times$  & $\times$  & $\times$ \\
GPU-free consumer deployment  & \checkmark & N/A        & $\times$  & $\times$  & $\sim$ \\
\bottomrule
\end{tabular}
\end{table}

% ============================================================
\section{Results and Discussion}
% ============================================================

\subsection{Evaluation Setup}

Pilot evaluation was conducted on a single consumer-grade
laptop---device identifier ``S''---running Windows~11 Pro
(64-bit, x64-based processor) without any discrete GPU or
hardware accelerator.
The complete hardware specification is reported in
Table~\ref{tab:hardware}.

\begin{table}[htbp]
\centering
\caption{Pilot Evaluation Hardware Specification}
\label{tab:hardware}
\begin{tabular}{ll}
\toprule
\textbf{Component}     & \textbf{Specification} \\
\midrule
Processor              & 12th Gen Intel Core i5-12450HX \\
Clock Speed            & 2.40\,GHz base (up to 4.40\,GHz boost) \\
Cores / Threads        & 8P\,+\,4E cores, 16 threads \\
Architecture           & Alder Lake-HX, x64 \\
Installed RAM          & 16.0\,GB (15.7\,GB usable) \\
Operating System       & Windows 11 Pro, 64-bit \\
Discrete GPU           & None (Intel UHD Graphics integrated) \\
Touch / Pen Input      & Not available \\
\bottomrule
\end{tabular}
\end{table}

The i5-12450HX is a laptop-class hybrid-architecture
processor combining eight performance cores and four
efficiency cores.
Its selection as the evaluation platform is deliberate: it
represents a mid-range consumer device containing no discrete
GPU, no tensor processing unit, and no hardware inference
accelerator of any kind, constituting a conservative and
reproducible baseline for assessing the practical feasibility
of fully local AI inference.
All experiments were executed within a Python~3.11 virtual
environment under standard background OS load; no CUDA
libraries were loaded at any stage.

The evaluation corpus comprised 200 manually authored journal
entries spanning an eight-week period, with 35 entry pairs
independently annotated for semantic relatedness by three
domain-expert human raters (inter-rater agreement
$\kappa = 0.81$, substantial).

\subsection{Performance Metrics}

Table~\ref{tab:perf} summarises the measured system
performance across all pipeline stages on the i5-12450HX
evaluation device.

\begin{table}[htbp]
\centering
\caption{System Performance Metrics (Pilot Evaluation ---
Intel Core i5-12450HX, 16\,GB RAM, No GPU)}
\label{tab:perf}
\begin{tabular}{lcc}
\toprule
\textbf{Metric}              & \textbf{Measured} & \textbf{Target} \\
\midrule
End-to-end query latency     & 1.04\,s           & $<1.2$\,s \\
RoBERTa embedding time       & 183\,ms           & $<250$\,ms \\
FAISS insert latency         & 769\,ms           & $<800$\,ms \\
FAISS top-5 retrieval time   & 9\,ms             & $<50$\,ms \\
Rasa NLU intent latency      & 71\,ms            & $<100$\,ms \\
Semantic link precision      & 80\%              & $>75\%$ \\
Baseline static precision    & 60\%              & --- \\
Peak RAM usage               & 5.1\,GB           & $<12$\,GB \\
D3.js timeline frame rate    & 60\,FPS           & 60\,FPS \\
\bottomrule
\end{tabular}
\end{table}

\subsection{Semantic Link Quality}

The Adaptive Min-P threshold $T_{\mathrm{adaptive}}$
(Eq.~\ref{eq:adaptive}) yielded a semantic link precision of
80\% against human annotation, a relative improvement of
33.3\% over the fixed $\tau = 0.78$ baseline (60\%).
The improvement is attributable to the threshold's sensitivity
to each user's personal vocabulary distribution: for the
evaluation author, whose writing exhibited high lexical
diversity, $T_{\mathrm{adaptive}} = 0.71$ (lower than the
static baseline), correctly admitting cross-domain
associations that the static threshold rejected.

\subsection{Cognitive Drift Validity}

Cognitive Drift scores computed over rolling 7-day windows
correlated positively ($r = 0.67$, $p < 0.01$) with
self-reported perceived stress ratings collected concurrently
via a five-point Likert scale, suggesting that the metric
captures a psychologically meaningful signal despite its
deliberately non-pathological formulation.
Formal validation against standardised psychometric
instruments is deferred to the longitudinal study outlined in
Section~VI.

\subsection{Discussion}

All primary latency targets were met on GPU-free consumer
hardware (Intel Core i5-12450HX, 16\,GB RAM), confirming the
architectural feasibility of the local-first constraint even
on mid-range laptop hardware without specialist inference
accelerators.
The 33.3\% precision gain from adaptive thresholding over
static baselines, validated against expert annotation,
provides quantitative justification for the design choice.
The Rumination Guard was not triggered during the eight-week
pilot (no participant exhibited three or more consecutive
high-similarity negative entries), indicating appropriate
specificity.
Limitations include the relatively small evaluation corpus
($n = 200$), the single-user pilot design, and the absence
of a longitudinal clinical control condition.

% ============================================================
\section{Conclusion and Future Work}
% ============================================================

This paper has presented \CC{}, a fully local,
privacy-preserving spatial journaling framework that
demonstrates the practical feasibility of combining dense
semantic retrieval (FAISS, RoBERTa), on-device intent
classification (Rasa NLU), adaptive graph construction, and
temporally structured visualisation (Rigid Narrative Timeline)
on consumer-grade hardware without cloud dependency.

The system closes four documented gaps in existing AI journaling
tools: cloud-mandatory inference, static semantic thresholding,
temporally uninformed graph topology, and unvalidated clinical
framing.
Pilot evaluation confirms that all primary latency targets are
achievable without GPU acceleration, and that the Adaptive
Min-P thresholding protocol delivers a 33.3\% improvement in
semantic link precision relative to the static-threshold
baseline.

Four directions are identified for future work:

\begin{enumerate}
  \item \textbf{Longitudinal clinical validation.}
    A pre-registered study involving a cohort of
    $n \geq 100$ participants over 12 weeks, with
    standardised psychometric measures (PHQ-9, PSS-10), is
    required to assess the clinical correlates of the
    Cognitive Drift metric and Rumination Guard.

  \item \textbf{Multimodal entry support.}
    Extension of the embedding pipeline to handle voice memos
    and handwritten image inputs via local
    Whisper~\cite{radford2023whisper} and Tesseract OCR would
    broaden the system's practical utility without violating
    the local-first constraint.

  \item \textbf{Federated personalisation.}
    Privacy-preserving federated learning over encrypted local
    model checkpoints could enable community-level threshold
    calibration and intent model improvement without
    centralising any user data~\cite{li2020federated}.

  \item \textbf{Scalability under large corpora.}
    The current \texttt{IndexFlatIP} configuration is
    appropriate for corpora up to approximately 10,000
    entries.
    Evaluation of \texttt{IndexHNSW} and
    \texttt{IndexIVFPQ} configurations is warranted for users
    with decade-long journaling histories, where
    $N > 50{,}000$ entries are plausible.
\end{enumerate}

% ============================================================
\section*{Acknowledgment}
% ============================================================

The authors thank the volunteer participants of the pilot
evaluation and the anonymous reviewers whose constructive
feedback improved this manuscript.

% ============================================================
% REFERENCES
% ============================================================
\bibliographystyle{IEEEtran}

\begin{thebibliography}{99}

\bibitem{pennebaker1997}
J.~W. Pennebaker and J.~D. Seagal,
``Forming a story: The health benefits of narrative,''
\textit{J. Clin. Psychol.},
vol.~55, no.~10, pp.~1243--1254, 1997.

\bibitem{smyth1998}
J.~M. Smyth,
``Written emotional expression: Effect sizes, outcome types,
and moderating variables,''
\textit{J. Consult. Clin. Psychol.},
vol.~66, no.~1, pp.~174--184, 1998.

\bibitem{xu2023mindfuldiary}
T.~Kim, S.~Bae, H.~A. Kim, S.-W. Lee, H.~Hong,
C.~Yang, and Y.-H. Kim,
``MindfulDiary: Harnessing large language model to support
psychiatric patients' journaling,''
in \textit{Proc. 2024 CHI Conf. Human Factors Comput. Syst.}
(CHI~'24), New York, NY, USA: ACM, Article 701,
pp.~1--20, 2024.
\url{https://doi.org/10.1145/3613904.3642937}

\bibitem{feretzakis2024}
G.~Feretzakis, K.~Papaspyridis, A.~Gkoulalas-Divanis,
and V.~S. Verykios,
``Privacy-preserving techniques in generative AI and large
language models: A narrative review,''
\textit{Information}, vol.~15, p.~697, 2024.

\bibitem{tauscher2023}
J.~S. Tauscher \textit{et al.},
``Automated detection of cognitive distortions in text
exchanges between clinicians and people with serious mental
illness,''
\textit{Psychiatric Services},
vol.~74, no.~4, pp.~407--410, 2023.

\bibitem{sage2025}
A.~Sage, J.~Keppens, and H.~Yannakoudakis,
``A survey of cognitive distortion detection and
classification in NLP,''
in \textit{Findings of ACL: EMNLP 2025},
pp.~14884--14899, 2025.

\bibitem{mindscape}
S.~Nepal, A.~Pillai, W.~Campbell, T.~Massachi,
M.~V. Heinz, A.~Kunwar, E.~S. Choi, X.~Xu,
J.~Kuc, J.~F. Huckins, J.~Holden, S.~M. Preum,
C.~Depp, N.~Jacobson, M.~P. Czerwinski,
E.~Granholm, and A.~T. Campbell,
``MindScape study: Integrating LLM and behavioral sensing
for personalized AI-driven journaling experiences,''
\textit{Proc. ACM Interact. Mob. Wearable Ubiquitous
Technol.}, vol.~8, no.~4, Article 186,
pp.~1--44, Dec.~2024.
\url{https://doi.org/10.1145/3699761}

\bibitem{cjournal2024}
N.~Elsharawi and A.~El~Bolock,
``C-Journal: A journaling application for detecting and
classifying cognitive distortions using deep-learning based
on a crowd-sourced dataset,''
in \textit{Proc. 2024 Joint Int. Conf. Comput. Linguist.,
Lang. Resources Eval.} (LREC-COLING~2024),
Torino, Italia: ELRA and ICCL,
pp.~3224--3234, 2024.

\bibitem{lewis2020rag}
P.~Lewis \textit{et al.},
``Retrieval-augmented generation for knowledge-intensive
NLP tasks,''
in \textit{Proc. NeurIPS}, 2020, pp.~9459--9474.

\bibitem{karpukhin2020dpr}
V.~Karpukhin \textit{et al.},
``Dense passage retrieval for open-domain question
answering,''
in \textit{Proc. EMNLP}, 2020, pp.~6769--6781.

\bibitem{johnson2021faiss}
J.~Johnson, M.~Douze, and H.~J\'{e}gou,
``Billion-scale similarity search with GPUs,''
\textit{IEEE Trans. Big Data},
vol.~7, no.~3, pp.~535--547, 2021.

\bibitem{dettmers2022llmint8}
T.~Dettmers, M.~Lewis, Y.~Belkada, and L.~Zettlemoyer,
``LLM.int8(): 8-bit matrix multiplication for transformers
at scale,''
in \textit{Proc. NeurIPS}, 2022.

\bibitem{liu2019roberta}
Y.~Liu \textit{et al.},
``RoBERTa: A robustly optimized BERT pretraining approach,''
\textit{arXiv preprint arXiv:1907.11692}, 2019.

\bibitem{bunk2020diet}
T.~Bunk, D.~Varshneya, V.~Vlasov, and A.~Nichol,
``DIET: Lightweight language understanding for dialogue
systems,''
\textit{arXiv preprint arXiv:2004.09936}, 2020.

\bibitem{pennebaker2003liwc}
J.~W. Pennebaker, M.~R. Mehl, and K.~G. Niederhoffer,
``Psychological aspects of natural language use: Our words,
our selves,''
\textit{Annu. Rev. Psychol.},
vol.~54, pp.~547--577, 2003.

\bibitem{welford1962}
B.~P. Welford,
``Note on a method for calculating corrected sums of squares
and products,''
\textit{Technometrics},
vol.~4, no.~3, pp.~419--420, 1962.

\bibitem{he2021deberta}
P.~He, X.~Liu, J.~Gao, and W.~Chen,
``DeBERTa: Decoding-enhanced BERT with disentangled
attention,''
in \textit{Proc. ICLR}, 2021.

\bibitem{aho1972}
A.~V. Aho, M.~R. Garey, and J.~D. Ullman,
``The transitive reduction of a directed graph,''
\textit{SIAM J. Comput.},
vol.~1, no.~2, pp.~131--137, 1972.

\bibitem{amershi2014}
S.~Amershi, M.~Cakmak, W.~B. Knox, and T.~Kulesza,
``Power to the people: The role of humans in interactive
machine learning,''
\textit{AI Mag.}, vol.~35, no.~4, pp.~105--120, 2014.

\bibitem{horvitz1999}
E.~Horvitz,
``Principles of mixed-initiative user interfaces,''
in \textit{Proc. ACM CHI}, 1999, pp.~159--166.

\bibitem{greenberger1995}
D.~Greenberger and C.~A. Padesky,
\textit{Mind Over Mood: Change How You Feel by Changing
the Way You Think}.
New York, NY: Guilford Press, 1995.

\bibitem{radford2023whisper}
A.~Radford \textit{et al.},
``Robust speech recognition via large-scale weak
supervision,''
in \textit{Proc. ICML}, 2023.

\bibitem{li2020federated}
T.~Li \textit{et al.},
``Federated learning: Challenges, methods, and future
directions,''
\textit{IEEE Signal Process. Mag.},
vol.~37, no.~3, pp.~50--60, 2020.

\end{thebibliography}

\end{document}
